\section*{Exercises}
\addcontentsline{toc}{section}{Exercises}

The solutions to the exercises that do not ask for code are in
\hyperref[ch:losningar]{appendix~A}. The code is published in the course repository after the
lecture.

\exercise{Build a frame}{Calculation}
\label{ex:1:1}
A node with address \code{0x17} sends a \code{StatusRequest} to a temperature sensor with address
\code{0x25}. The sequence number is \code{0x7F05}.
\begin{enumerate}[a)]
  \item Write out the whole frame byte by byte, in hexadecimal.
  \item The sensor answers with a \code{StatusResponse} containing the temperature \code{0x16}
    (one byte). Write out that frame as well.
  \item How many bytes is the shortest possible frame in this protocol?
\end{enumerate}

\exercise{The fields and their purpose}{Reasoning}
\label{ex:1:2}
For each of the fields SOF, LEN, TYPE, DST, SRC, SEQ and CHK: describe in one sentence what
breaks if the field is removed.

Two of them can be removed without the protocol ceasing to work in a system with exactly two
nodes. Which, and why?

\exercise{Endianness}{Reasoning}
\label{ex:1:3}
A group serialises the sequence number by copying memory directly:

\begin{cppcode}
std::memcpy(&buf[Offset::Seq], &frame.seqNr, sizeof(frame.seqNr));
\end{cppcode}

\begin{enumerate}[a)]
  \item What ends up in the buffer on an x86 machine if \code{seqNr} is \code{0x7F05}?
  \item The code nevertheless works flawlessly in their own tests. Why?
  \item When does it break?
\end{enumerate}

\exercise{The limits of the checksum}{Reasoning}
\label{ex:1:4}
The course's checksum is a sum of bytes.
\begin{enumerate}[a)]
  \item A disturbance swaps two data bytes. Is it detected?
  \item A disturbance increases one byte by 1 and decreases another by 1. Is it detected?
  \item Give an example of a change to the frame above that the checksum does \emph{not} detect.
\end{enumerate}
